\begin{enumerate}
\item Dans le repère $\big(\mathrm{A}~;~\vect{i},\vect{j},\vect{k}\big)$, on a les coordonnées suivantes pour les sommets du prisme droit :

	\hfill~A$(0~;~0~;~0)$ \hfill B$ (4~;~0~;~0) $ \hfill C$(4~;~4~;~0)$ \hfill D$(0~;~4~;~0)$\hfill~

	\hfill~E$(0~;~0~;~8)$ \hfill F$ (4~;~0~;~4) $ \hfill G$(4~;~4~;~4)$ \hfill H$(0~;~4~;~8)$\hfill~

	I étant le milieu de [EF], on a I$\left(\dfrac{x_\mathrm{E}+x_\mathrm{F}}{2}~;~\dfrac{y_\mathrm{E}+y_\mathrm{F}}{2}~;~\dfrac{z_\mathrm{E}+z_\mathrm{F}}{2}\right)$
	, soit I$(2~;~0~;~6)$.

	J étant le milieu de [AE], on a de même : J$(0~;~0~;~4)$.

\item 
	\begin{enumerate}
		\item Si le plan est nommé (IGJ), cela signifie que les trois points I, G et J définissent le plan, et donc sont non alignés.

		On a : $\vect{\mathrm{IG}}\begin{pmatrix} x_\mathrm{G}-x_\mathrm{I}\\y_\mathrm{G}-y_\mathrm{I}\\z_\mathrm{G}-z_\mathrm{I} \end{pmatrix} = \begin{pmatrix} 2\\4\\-2\end{pmatrix}$ \quad et de même : $\vect{\mathrm{IJ}} \begin{pmatrix} -2\\0\\-2\end{pmatrix}$ .

		Comme le repère est orthonormé, on peut calculer les produits scalaires à l'aide des coordonnées :

\begin{itemize}
\item $ \vect{n}\cdot \vect{\mathrm{IG}} = -1\times 2 + 1 \times 4 + 1 \times (-2) = -2 + 4-2 = 0$ : $\vect{n}$ est orthogonal à $ \vect{\mathrm{IG}} $.
\item $ \vect{n}\cdot \vect{\mathrm{IJ}} = -1\times (-2) + 1 \times 0 + 1\times (-2) = 2 + 0-2 = 0$ : $\vect{n}$ est orthogonal à $ \vect{\mathrm{IJ}} $.
\end{itemize}

Ainsi, $\vect{n}$ est orthogonal à deux vecteurs non colinéaires du plan (IGJ) : c'est donc un vecteur normal au plan.

	\item Une équation cartésienne d'un plan dont $ \vect{n} $ est un vecteur normal est de la forme : \linebreak$-1\times x + 1 \times y + 1\times z + d =0$, soit $-x+y+z+d=0$, où $d$ est un réel quelconque.

Comme G est un point du plan (IGJ), on en déduit que la constante $d$ dans ce cas doit être telle que :

$\begin{aligned} -x_\mathrm{G} +y_\mathrm{G}+z_\mathrm{G} +d=0 &\iff -4+4+4+d = 0\\&\iff d = -4\\ \end{aligned}$

Une équation de (IGJ) est donc : $-x + y + z - 4 = 0$.
	\end{enumerate}
\item Si $d$ est perpendiculaire à (IGJ), alors elle est dirigée par $\vect{n}$, comme elle passe par H, elle admet comme représentation paramétrique : $ \left\{\begin{array}{l} x=x_\mathrm{H} + tx_{\vect{n}}\\y=y_\mathrm{H} + ty_{\vect{n}}\\z=z_\mathrm{H} + tz_{\vect{n}} \end{array}\right. \quad t \in \R $, soit ici : $ \left\{\begin{array}{l} x=-t\\y=4 + t\\z=8+t \end{array}\right. \quad t \in \R $

\item Si L est le projeté orthogonal de H sur (IGJ), cela veut dire que la droite (HL) est orthogonale au plan (et passe par H), et donc que la droite (HL) est la droite $d$. Comme L est un point du plan, c'est donc le seul point de $d$ sur le plan.

Cherchons le paramètre $t$ tel qu'un point $M_t$ de paramètre $t$ dans la représentation de $d$ soit un point de (IGJ) :

$\begin{aligned} M_t \in (\mathrm{IGJ}) &\iff -x_{M_t} + y_{M_t}+z_{M_t} - 4 = 0\\
&\iff -(-t) + (4+t) + (8 + t) - 4 = 0\\
&\iff 3t +8 = 0\\
&\iff t = \dfrac{-8}{3}\\ \end{aligned}$

L est donc $M_{\frac{-8}{3}}$ sur la droite $d$ : il a donc comme coordonnées L$\left(-\dfrac{-8}{3}~;~4 + \dfrac{-8}{3}~;~8 + \dfrac{-8}{3} \right)$.

Cela confirme L$\left(\dfrac{8}{3}~;~\dfrac{12-8}{3} ~;~\dfrac{24-8}{3}\right) $

Autrement dit, le point L est bien le point de coordonnées $ \left(\dfrac{8}{3}~;~\dfrac{4}{3}~;~\dfrac{16}{3}\right)  $.

\item Par définition, la distance d'un point à un plan est la distance entre le point et son projeté orthogonal sur le plan, donc on cherche HL. Comme on travaille dans un repère orthonormé :

$\mathrm{HL} = \sqrt{(x_\mathrm{L} - x_\mathrm{H})^2+(y_\mathrm{L} - y_\mathrm{H})^2+(z_\mathrm{L} - z_\mathrm{H})^2} = \sqrt{\left(\dfrac{8}{3}\right)^2+\left(\dfrac{8}{3}\right)^2+\left(\dfrac{8}{3}\right)^2 } = \sqrt{\dfrac{192}{9}}=\dfrac{8\sqrt{3}}{3}$.

La distance de H au plan (IGJ) est donc de $ \dfrac{8\sqrt{3}}{3}$.

\item Calculons : $\vect{\mathrm{IG}}\cdot\vect{\mathrm{IJ}} = 2\times (-2) + 4\times 0 +(-2)\times (-2) = -4 + 0 +4= 0$.

Les vecteurs $\vect{\mathrm{IG}}$ et $\vect{\mathrm{IJ}}$ sont donc orthogonaux (et non nuls) donc les droites qu'ils dirigent, (IG) et (IJ) sont orthogonales (et perpendiculaires, car sécantes en I) : le triangle IGJ est donc rectangle en I.

\item Pour calculer le volume, on choisira IGJ comme base (car le triangle étant rectangle, son aire est simple à calculer) et donc la hauteur correspondante est la distance du quatrième sommet (H) au plan (IGJ) (distance qui a été calculée à la question \textbf{5.}).

$\mathrm{IG} = \sqrt{2^2 + 4^2 + (-2)^2} = \sqrt{24}=2\sqrt{6}$

$ \mathrm{IJ} = \sqrt{(-2)^2 + 0 + 2^2} = \sqrt{8} = 2\sqrt{2}$

L'aire du triangle IGJ est donc : $\mathcal{A}_{\mathrm{IGJ}} = \dfrac{\mathrm{IG}\times \mathrm{IJ}}{2} = \dfrac{2\sqrt{6}\times 2\sqrt{2}}{2} = 4\sqrt{3}$.

Le volume du tétraèdre est donc : $V = \dfrac{1}{3} \times 4\sqrt{3} \times \dfrac{8\sqrt{3}}{3} = \dfrac{32}{3}$.

Le volume du tétraèdre est de $V = \dfrac{32}{3}$ (soit environ 10,7, au dixième près).
\end{enumerate}

\bigskip
